%%======== TEMPLATE INFORMES ESTILO FCFM ========
% Autor: Elisa Paz Kauffmann Figueroa
% Estudiante de Ingeniería Civil en Ciencias de la Computación
% FCFM U. de Chile
%%===============================================

\documentclass[14pt,letterpaper,hidelinks]{extarticle}
%\usepackage[ansinew]{inputenc}
\usepackage[utf8x]{inputenc}
%\usepackage[latin1]{inputenc}
\usepackage[spanish]{babel}
\usepackage[letterpaper,includeheadfoot, top=1cm, bottom=3.0cm, right=2.0cm, left=2.0cm]{geometry}
\renewcommand{\familydefault}{\rmdefault}

\usepackage[table,xcdraw]{xcolor}
\usepackage{svg}
%\usepackage{graphicx}
%\usepackage{color}
\usepackage{hyperref}
\usepackage{amssymb}
\usepackage{url}
%\usepackage{pdfpages}
\usepackage{fancyhdr}
\usepackage{hyperref}
\usepackage{subfig}
\usepackage{indentfirst}
\usepackage{titlesec}
\usepackage{booktabs}
\usepackage{multirow}
\usepackage{listings} %Codigo
\lstset{language=C, tabsize=4,framexleftmargin=5mm,breaklines=true}

\titleformat*{\section}{\large\bfseries}
\titleformat*{\subsection}{\normalsize\bfseries}

\begin{document}
\thispagestyle{empty}
\renewcommand*\listtablename{Índice de tablas}
\renewcommand{\contentsname}{\'Indice}
\renewcommand*{\refname}{}

%\begin{sf}
% --------------- ---------PORTADA --------------------------------------------
\newpage
\pagestyle{fancy}
\fancyhf{}
%-------------------- CABECERA ---------------------
\hbox{\includegraphics[scale=0.3]{img/fcfm_dcc_png.png} }
%------------------ TÍTULO -----------------------
\vspace*{4cm}
\begin{center}
\Huge  {Tarea 2 - Comerciales}\\
\vspace{1cm}
\Large {CC5204 - Búsqueda por contenido de Imágenes y Videos}\\


\vspace{1cm}
\end{center}
%----------------- NOMBRES ------------------------
\begin{flushright}
\begin{table}[h]
	\large
	\centering
	\begin{tabular}{ll}
		Alumnos: & Patricio Isbej\\
				& Elisa Kauffmann\\
		Profesor: & Juan Manuel Barrios\\
		Ayudante: & Sebastián Ferrada\\
		Fecha: & 2 de mayo de 2016
	\end{tabular}
\end{table}
\end{flushright}

% ·············· ENCABEZADO - PIE DE PAGINA ············
\newpage
\pagestyle{fancy}
\fancyhf{}

%Encabezado
%\fancyhead[L]{\rightmark}
\fancyhead[L]{\small \rm \textit{Sección \nouppercase\rightmark}} %Izquierda
\fancyhead[R]{\small \rm \textbf{\thepage}} %Derecha


%\fancyfoot[L]{\small \rm \textit{Pie de página - Izquierda}} %Izquierda
%\fancyfoot[R]{\small \rm \textit{Pie de página - Derecha}} %Derecha
%\fancyfoot[C]{\thepage} %Centro

\renewcommand{\sectionmark}[1]{\markright{\thesection.\ #1}}
\renewcommand{\headrulewidth}{0.5pt}
%\renewcommand{\footrulewidth}{0.5pt}

% =============== INDICE ===============
%
\tableofcontents
\listoffigures
\listoftables

% =============== SECCIONES Y SUBSECCIONES ===============
\newpage
\section{Introducción}
\section{Resumen}
\section{Hipótesis}
\section{Diseño e implementación}
Se implementa un detector de comerciales que señala cuántas veces fue emitido cada comercial durante las 4 horas de programación, determinando exactamente el inicio y
fin de cada aparición. Para ello se divide el trabajo en 3 pasos: 
\begin{enumerate}
\item Descripción del contenido (extracción de características)
\item Búsqueda por similitud (comparación de descriptores)
\item Detección (encontrar secuencias parecidas)
\end{enumerate}

Se usa el lenguaje \verb+Python+ utilizando las librerías \verb+OpenCV+ para procesar imágenes y videos, \verb+numpy+ para manejo de vectores y matrices y \verb+matplotlib+ para graficar.\\

El reporte de detecciones se guarda en un archivo \verb+.m3u+ que por cada aparición de un comercial, contiene 3 líneas con el siguiente formato: 
\begin{verbatim}
#EXTVLCOPT:start-time= [seconds]
#EXTVLCOPT:stop-time= [seconds]
../../base/mega-2014_04_25T22_00_07.mp4
\end{verbatim}
donde la primera y la segunda línea indican el momento en que empieza y termina el comercial respectivamente (en segundos), y la tercera línea indica el video a reproducir, en este caso es el programa de 4 horas.\\

\subsection{Descripción del contenido}
A rasgos generales, en esta etapa se utiliza el descriptor de \textbf{histogramas por zonas} para caracterizar los frames de los videos guardando el descriptor de cada video en un archivo \verb+.npy+, el cual será utilizado en el paso siguiente. Es importante mencionar que los frames de todos los videos son escalados al mismo tamaño ($720\times400$ px) para normalizar y facilitar la división en zonas. Además, previo al desarrollo de la tarea se verificó que todos los videos tuvieran el mismo \textit{framerate} para mantener consistencia ($\sim$ 29.97 fps).\\

En esta etapa, los parámetros que varían en cada configuración son:
\begin{itemize}
\item[-] frecuencia de muestreo de frames
\item[-] cantidad de zonas en las que se divide cada frame
\item[-] cantidad de bins de los histogramas 
\end{itemize}

En detalle, el programa ejecuta los siguientes pasos:
\begin{enumerate}
\item Para cada archivo de video, se toma un frame cada cierta frecuencia, lo escala a un tamaño de $720\times400$ px y lo transforma a escala de grises.
\item Se divide cada uno de estos frames en la cantidad de zonas especificada por el tipo de configuración, y usando la función \verb+calcHist+ de \verb+OpenCV+ se obtiene un histograma por cada zona en forma de un arreglo.
\item Se crea el descriptor completo del frame: un arreglo que es la concatenación de los vectores de los histogramas obtenidos en el paso anterior.
\item Se crea el descriptor del video completo construyendo un arreglo \verb+numpy+ de todos los arreglos que describen cada frame del video.
\end{enumerate}

El código se encuentra en el archivo \verb+describe.py+

\subsection{Búsqueda por similitud}
Durante esta etapa, usando los descriptores creados anteriormente, por cada frame del video de 4 horas se encuentra el frame del comercial que más se le parece. Para ello se comparan todos los vectores descriptores de cada comercial contra cada frame del video de 4 horas, utilizando distancia Manhattan. Para facilitar el procesamiento se crea un diccionario que guarda, por cada comercial, su título, un número identificador y la cantidad de frames procesados. Un ejemplo de uno de los diccionarios creados se puede observar en el cuadro~\ref{tab:dict}.\\

\begin{table}[]
\centering
\begin{tabular}{@{}ccc@{}}
\toprule
Id & Nombre                      & Frames \\ \midrule
0  & mr big (2).npy              & 95     \\
1  & scotiabank.npy              & 76     \\
2  & ripley dias r.npy           & 63     \\
3  & sodimac homecenter.npy      & 62     \\
4  & donnasept.npy               & 61     \\
5  & dove desodorante hombre.npy & 45     \\
6  & entel 4g corto.npy          & 81     \\
7  & johnson mundial.npy         & 60     \\
8  & limon soda.npy              & 96     \\
9  & bilz y pap.npy              & 105    \\ \bottomrule
\end{tabular}
\captionsetup{justification=centering,margin=2cm}
\caption{Bosquejo del diccionario creado para la configuración de 1 descriptor cada 10 frames. \label{tab:dict}}
\end{table}

Específicamente, el programa realiza los siguientes pasos:
\begin{enumerate}
\item Se crea un diccionario que mapea los títulos de los comerciales a un identificador y a la cantidad de frames obtenidos en la etapa anterior.
\item Se cargan en memoria los descriptores de video de todos los 10 comerciales y usando el diccionario se construye una matriz de dimensiones $10\times 3$. La matriz contiene, para cada comercial, una fila con su identificador, su título y su descriptor. (Ver cuadro~\ref{tab:descsCom}). 
\item Se carga el descriptor del programa de 4 horas y se compara cada frame contra todos los frames de los 10 comerciales. Para medir la similitud entre vectores descriptores se utiliza distancia Manhattan (utilizando la función \verb+norm+ de \verb+OpenCV+ con el parámetro \verb+NORM_L1+). Así se va construyendo un vector de ``calces'' con tantas filas como frames del programa de 4 horas procesados. Cada fila contiene el índice del frame y el identificador del comercial que más se parece al frame del programa.
\item El vector de ``calces'' se guarda en un archivo \verb+.npy+ para ser usado en el paso siguiente   
\end{enumerate} 
\begin{table}[]
\centering
\begin{tabular}{@{}ccc@{}}
\toprule
com\_id & nombre archivo & descriptor comercial \\ \midrule
0 & mr big (2).npy & {[}{[}{]},{[}{]},...,{[}{]},{[}{]}{]} \\
1 & scotiabank.npy & {[}{[}{]},{[}{]},...,{[}{]},{[}{]}{]} \\
2 & ripley dias r.npy & {[}{[}{]},{[}{]},...,{[}{]},{[}{]}{]} \\
3 & sodimac homecenter.npy & {[}{[}{]},{[}{]},...,{[}{]},{[}{]}{]} \\
4 & donnasept.npy & {[}{[}{]},{[}{]},...,{[}{]},{[}{]}{]} \\
5 & dove desodorante hombre.npy & {[}{[}{]},{[}{]},...,{[}{]},{[}{]}{]} \\
6 & entel 4g corto.npy & {[}{[}{]},{[}{]},...,{[}{]},{[}{]}{]} \\
7 & johnson mundial.npy & {[}{[}{]},{[}{]},...,{[}{]},{[}{]}{]} \\
8 & limon soda.npy & {[}{[}{]},{[}{]},...,{[}{]},{[}{]}{]} \\
9 & bilz y pap.npy & {[}{[}{]},{[}{]},...,{[}{]},{[}{]}{]} \\ \bottomrule
\end{tabular}
\captionsetup{justification=centering,margin=2cm}
\caption{Bosquejo de la matriz de identificadores y descriptores de los comerciales \label{tab:descsCom}}
\end{table}

El código se encuentra en el archivo \verb+compare.py+

\subsection{Detección}
En esta etapa se detecta la aparición de los comerciales en el video de 4 horas haciendo uso del diccionario y el vector de ``calces'' creados en las etapas anteriores. La idea general es usar máscaras, una por cada comercial, para ir recorriendo el vector de ``calces'' y comparar la similitud entre segmentos de éste con las máscaras usando distancia de Hamming. La idea se ilustra en la figura~\ref{my-label}.

La distancia de Hamming es calculada usando una función específica, implementada para este caso en particular.\\

Detalladamente, la implementación lleva a cabo los siguientes pasos:
\begin{enumerate}
\item Se cargan el diccionario y el vector de ``calces'' en memoria a partir de los archivos generados en la etapa anterior.
\item Para cada comercial en el diccionario se crea una máscara: una matriz con tantas filas como la cantidad de frames del comercial y 2 columnas. En la primera columna se repite para todas las filas el identificador del comercial y en la segunda columna se enumeran todos sus frames. Un ejemplo de máscara se puede ver en la figura~\ref{tab:mask}.
\item Para cada línea del vector de ``calces'' se alínea la máscara del comercial indicado en la posición del frame indicado. Se calcula la similitud entre la máscara y la sección del vector de calces usando distancia de Hamming
\end{enumerate}

\begin{table}[]
\centering
\begin{tabular}{@{}cc@{}}
\toprule
Id comercial & frame \\ \midrule
2 & 0 \\
2 & 1 \\
2 & 2 \\
2 & 3 \\
2 & 4 \\
2 & 5 \\
2 & 6 \\
2 & 7 \\
2 & 8 \\
2 & 9 \\
2 & 10 \\
2 & 11 \\
2 & 12 \\
2 & 13 \\
2 & 14 \\ \bottomrule
\end{tabular}
\captionsetup{justification=centering,margin=2cm}
\caption{Bosquejo de la máscara de un comercial creada para la configuración de 1 descriptor cada 30 frames\label{tab:mask}}
\end{table}

\section{Experimentos y resultados}
\begin{table}[]
\centering
\caption{My caption}
\label{my-label}
\begin{tabular}{c|c|c|c|c|c|c|c|}
\cline{2-8}
 & Frecuencia & \multicolumn{2}{c|}{30} & \multicolumn{2}{c|}{20} & \multicolumn{2}{c|}{10} \\ \hline
\multicolumn{1}{|c|}{Division} & Bins$\backslash$Tolerancia & 50\% & 80\% & 50\% & 80\% & 50\% & 80\% \\ \hline
\multicolumn{1}{|c|}{\multirow{3}{*}{1x1}} & 16 & 0.75 & 0.98 & 0.81 & 1 & 0.84 & 1 \\ \cline{2-8} 
\multicolumn{1}{|c|}{} & 32 & 0.75 & 1 & 0.78 & 1 & 0.89 & 1 \\ \cline{2-8} 
\multicolumn{1}{|c|}{} & 64 & 0.78 & 1 & 0.81 & 1 & 0.89 & 1 \\ \hline
\multicolumn{1}{|c|}{\multirow{3}{*}{2x2}} & 16 & 0.86 & 1 & 0.94 & 1 & 0.94 & 1 \\ \cline{2-8} 
\multicolumn{1}{|c|}{} & 32 & 0.89 & 1 & 0.86 & 1 & 0.94 & 1 \\ \cline{2-8} 
\multicolumn{1}{|c|}{} & 64 & 0.89 & 0.98 & 0.91 & 1 & 0.94 & 1 \\ \hline
\multicolumn{1}{|c|}{\multirow{3}{*}{4x4}} & 16 & 0.91 & 0.96 & 0.89 & 1 & 0.96 & 1 \\ \cline{2-8} 
\multicolumn{1}{|c|}{} & 32 & 0.91 & 0.96 & 0.86 & 1 & 0.94 & 1 \\ \cline{2-8} 
\multicolumn{1}{|c|}{} & 64 & 0.91 & 0.94 & 0.89 & 1 & 0.96 & 1 \\ \hline
\end{tabular}
\end{table}

\begin{table}[]
\centering
\caption{My caption}
\label{my-label}
\begin{tabular}{c|
>{\columncolor[HTML]{C0C0C0}}c |c|c|c|c|c|c|}
\cline{2-8}
                                                                                           & \cellcolor[HTML]{9B9B9B}{\color[HTML]{000000} Frecuencia} & \multicolumn{2}{c|}{\cellcolor[HTML]{9B9B9B}{\color[HTML]{000000} 30}} & \multicolumn{2}{c|}{\cellcolor[HTML]{9B9B9B}{\color[HTML]{000000} 20}} & \multicolumn{2}{c|}{\cellcolor[HTML]{9B9B9B}{\color[HTML]{000000} 10}} \\ \hline
\multicolumn{1}{|c|}{\cellcolor[HTML]{9B9B9B}{\color[HTML]{000000} Division}}              & Bins$\backslash$Tolerancia                                           & \cellcolor[HTML]{C0C0C0}50\%       & \cellcolor[HTML]{C0C0C0}80\%      & \cellcolor[HTML]{C0C0C0}50\%       & \cellcolor[HTML]{C0C0C0}80\%      & \cellcolor[HTML]{C0C0C0}50\%       & \cellcolor[HTML]{C0C0C0}80\%      \\ \hline
\multicolumn{1}{|c|}{\cellcolor[HTML]{9B9B9B}{\color[HTML]{000000} }}                      & {\color[HTML]{000000} 16}                                 & 15/0/10                            & 24/0/1                            & 17/0/8                             & 25/0/0                            & 18/0/7                             & 25/0/0                            \\ \cline{2-8} 
\multicolumn{1}{|c|}{\cellcolor[HTML]{9B9B9B}{\color[HTML]{000000} }}                      & {\color[HTML]{000000} 32}                                 & 15/0/10                            & 25/0/0                            & 16/0/9                             & 25/0/0                            & 20/0/5                             & 25/0/0                            \\ \cline{2-8} 
\multicolumn{1}{|c|}{\multirow{-3}{*}{\cellcolor[HTML]{9B9B9B}{\color[HTML]{000000} 1x1}}} & {\color[HTML]{000000} 64}                                 & 16/0/9                             & 25/0/0                            & 20/0/5                             & 25/0/0                            & 20/0/5                             & 25/0/0                            \\ \hline
\multicolumn{1}{|c|}{\cellcolor[HTML]{9B9B9B}{\color[HTML]{000000} }}                      & {\color[HTML]{000000} 16}                                 & 19/0/6                             & 25/0/0                            & 22/0/3                             & 25/0/0                            & 22/0/3                             & 25/0/0                            \\ \cline{2-8} 
\multicolumn{1}{|c|}{\cellcolor[HTML]{9B9B9B}{\color[HTML]{000000} }}                      & {\color[HTML]{000000} 32}                                 & 20/0/5                             & 25/0/0                            & 19/0/6                             & 25/0/0                            & 22/0/3                             & 25/0/0                            \\ \cline{2-8} 
\multicolumn{1}{|c|}{\multirow{-3}{*}{\cellcolor[HTML]{9B9B9B}{\color[HTML]{000000} 2x2}}} & {\color[HTML]{000000} 64}                                 & 20/0/5                             & 25/1/0                            & 21/0/4                             & 25/0/0                            & 22/0/3                             & 25/0/0                            \\ \hline
\multicolumn{1}{|c|}{\cellcolor[HTML]{9B9B9B}{\color[HTML]{000000} }}                      & {\color[HTML]{000000} 16}                                 & 21/0/4                             & 25/2/0                            & 20/0/5                             & 25/0/0                            & 23/0/2                             & 25/0/0                            \\ \cline{2-8}
\multicolumn{1}{|c|}{\cellcolor[HTML]{9B9B9B}{\color[HTML]{000000} }}                      & {\color[HTML]{000000} 32}                                 & 21/0/4                             & 25/2/0                            & 19/0/6                             & 25/0/0                            & 22/0/3                             & 25/0/0                            \\ \cline{2-8} 
\multicolumn{1}{|c|}{\multirow{-3}{*}{\cellcolor[HTML]{9B9B9B}{\color[HTML]{000000} 4x4}}} & {\color[HTML]{000000} 64}                                 & 21/0/4                             & 25/3/0                            & 20/0/5                             & 25/0/0                            & 23/0/2                             & 25/0/0                            \\ \hline
\end{tabular}
\end{table}
\section{Análisis y conclusiones}
\section{Anexos}
\begin{table}[]
\centering
\caption{My caption}
\label{my-label}
\begin{tabular}{|cc|l|l|l|l|l|l|}
\hline
\multicolumn{1}{|l}{} & \multicolumn{1}{r|}{Frecuencia} & \multicolumn{2}{c|}{30}                               & \multicolumn{2}{c|}{20}                               & \multicolumn{2}{c|}{10}                               \\ \cline{3-8} 
\multicolumn{1}{|l}{Division} & \multicolumn{1}{l|}{Bins$\backslash$Tolerancia} & \multicolumn{1}{c|}{50\%} & \multicolumn{1}{c|}{80\%} & \multicolumn{1}{c|}{50\%} & \multicolumn{1}{c|}{80\%} & \multicolumn{1}{c|}{50\%} & \multicolumn{1}{c|}{80\%} \\ \hline
\multicolumn{1}{|c|}{\multirow{3}{*}{1x1}} & 16                                   & 15/0/10                   & 24/0/1                    & 17/0/8                    & 25/0/0                    & 18/0/7                    & 25/0/0                    \\ \cline{2-8} 
\multicolumn{1}{|c|}{}                     & 32                                   & 15/0/10                   & 25/0/0                    & 16/0/9                    & 25/0/0                    & 20/0/5                    & 25/0/0                    \\ \cline{2-8} 
\multicolumn{1}{|c|}{}                     & 64                                   & 16/0/9                    & 25/0/0                    & 20/0/5                    & 25/0/0                    & 20/0/5                    & 25/0/0                    \\ \hline
\multicolumn{1}{|c|}{\multirow{3}{*}{2x2}} & 16                                   & 19/0/6                    & 25/0/0                    & 22/0/3                    & 25/0/0                    & 22/0/3                    & 25/0/0                    \\ \cline{2-8} 
\multicolumn{1}{|c|}{}                     & 32                                   & 20/0/5                    & 25/0/0                    & 19/0/6                    & 25/0/0                    & 22/0/3                    & 25/0/0                    \\ \cline{2-8} 
\multicolumn{1}{|c|}{}                     & 64                                   & 20/0/5                    & 25/1/0                    & 21/0/4                    & 25/0/0                    & 22/0/3                    & 25/0/0                    \\ \hline
\multicolumn{1}{|c|}{\multirow{3}{*}{4x4}} & 16                                   & 21/0/4                    & 25/2/0                    & 20/0/5                    & 25/0/0                    & 23/0/2                    & 25/0/0                    \\ \cline{2-8} 
\multicolumn{1}{|c|}{}                     & 32                                   & 21/0/4                    & 25/2/0                    & 19/0/6                    & 25/0/0                    & 22/0/3                    & 25/0/0                    \\ \cline{2-8} 
\multicolumn{1}{|c|}{}                     & 64                                   & 21/0/4                    & 25/3/0                    & 20/0/5                    & 25/0/0                    & 23/0/2                    & 25/0/0                    \\ \hline
\end{tabular}
\end{table}

\begin{table}[]
\centering
\caption{My caption}
\label{my-label}
\begin{tabular}{ll}
2 & 0  \\
2 & 1  \\
2 & 2  \\
2 & 3  \\
2 & 4  \\
2 & 5  \\
2 & 6  \\
2 & 7  \\
2 & 8  \\
2 & 9  \\
2 & 10 \\
2 & 11 \\
2 & 12 \\
2 & 13 \\
2 & 14
\end{tabular}
\end{table}
\begin{table}[]
\centering
\caption{My caption}
\label{my-label}
\begin{tabular}{ll}
8 & 9  \\
4 & 16 \\
5 & 29 \\
5 & 29 \\
5 & 29 \\
4 & 16 \\
4 & 16 \\
2 & 1  \\
2 & 5  \\
2 & 3  \\
2 & 3  \\
2 & 5  \\
2 & 5  \\
2 & 7  \\
2 & 7  \\
2 & 8  \\
2 & 9  \\
2 & 11 \\
2 & 11 \\
2 & 12 \\
2 & 14 \\
2 & 14 \\
4 & 6  \\
4 & 6  \\
4 & 6  \\
4 & 6  \\
6 & 23 \\
4 & 15 \\
4 & 7  \\
5 & 25 \\
2 & 0  \\
4 & 3 
\end{tabular}
\end{table}

alineamiento\\

máscara\\

frame actual\\

calce\\

tamaño máscara\\
%%==================== IMAGENES =====================
%% ················ IMAGEN DOBLE ·················
%\begin{figure}[ht!] \centering
%\subfloat[Logo UChile]{\includegraphics[scale=0.2]{img/escudoU.pdf}}
%\hspace{1cm} %Espacio horizontal
%\subfloat[Logo FCFM]{\includegraphics[scale=0.45]{img/fcfm.png}}
%\caption{Ejemplo de imagen doble}\label{img1}
%\end{figure}
%%··········································
%
%
%A continuación la figura \ref{img2} presenta otra forma de agregar imágenes
%
%% ················ IMAGENES SIMPLES·················
%\begin{figure}[ht!]
%\centering \includegraphics[scale=0.2]{img/escudoU.pdf}
%\caption{Escudo de la Universidad de Chile} \label{img2}
%\end{figure}

%%--------------------------
%
%\begin{figure}[ht!]
%\centering 
%\captionsetup{justification=centering,margin=2cm}
%\includegraphics[scale=0.2]{img/fotos_alma/bunkers.JPG}
%\caption{Burros salvajes junto a los bunkers del campamento donde aloja el personal de OSF} 
%\label{campamento}
%\end{figure}  

%\begin{figure}[ht!]
%\centering \includegraphics[scale=0.2]{img/fotos_alma/cancha_casino.JPG}
%\caption{Multicancha y casino de OSF} 
%\label{cancha}
%\end{figure}  

%----------------------------
%\begin{figure}[ht!]
%\centering
%\hspace*{-2cm}
%\captionsetup{justification=centering,margin=2cm}
%\includegraphics[scale=0.3]{img/figure_2.pdf}
%\caption{Ejemplo de visualización de un perfil de temperatura usando datos de agosto de 2010} 
%\label{surftemp_fig}
%\end{figure}
%\begin{figure}
%\centering
%\hspace*{-2cm}
%\captionsetup{justification=centering,margin=2cm}
%\includegraphics[scale=0.3]{img/figure_3.pdf}
%\caption{Ejemplo de visualización de un perfil de temperatura usando datos de agosto de 2010} 
%\label{intliq_fig}
%\end{figure}
%%··········································



%%============== TABLAS ===============
%\begin{table}[h]
% \centering
% \caption{Headers, identificadores y descripción}
% \begin{tabular}{|| c | c | p{7	cm}||} 
% \hline
% Identificador de header & Tipo de medición & Descripción \\ [0.5ex] 
% \hline\hline
% 10 & - & No se usa (tampoco se describe en el manual)\\
% \hline
% 80 & - & No se usa (tampoco se describe en el manual) \\
% \hline
% 100 & 101 & Tipo de medición y título de los 4 tipos de perfiles\\ 
% \hline
% 200 & 201 & Header para mediciones meteorológicas a nivel de superficie\\
% \hline
% 300 & 301 & Header para mediciones escalares e integradas\\
% \hline
% 400 & 401, 402, 403, 404 & Ángulo de observación y arreglo de valores de 			altura, la varible independientele (58 valores, de 0 a 10 km), para todos 		los perfiles.\\
% \hline
% \end{tabular}
%\end{table}
%\newpage
%
%Las siguientes cuatro filas contienen el tipo de medición ``101", y especifican el tipo de medición para los cuatro tipos de perfiles.
%
%\begin{table}[h]
% \centering
% \caption{Tipo de perfil y datos}
% \begin{tabular}{||c|c|c|c|c||} 
% \hline
% N° de medición & Fecha/Tiempo & 100 &	Tipo de medición & Título \\ [0.5ex] 
% \hline\hline
% 1	& 10/07/14 22:07:14 & 101 & 401 & Temperature(K) \\
% \hline
% 2	& 10/07/14 22:07:14 & 101 & 402 & Vapor Density $(g/m^3)$ \\
% \hline
% 3	& 10/07/14 22:07:14 & 101 & 403 & Liquid Density $(g/m^3)$ \\
% \hline
% 4	& 10/07/14 22:07:14 & 101 & 404 & Relative Humidity (\%) \\
% \hline
% \end{tabular}
%\end{table}
%
%El resto de las filas son mediciones (de tipo 201, 301, 401, 402, 403 y 404). \textbf{Estas son las medidas que van a hacer almacenadas en la base de datos}
%\begin{table}[h]
% \centering
% \caption{Ejemplo de mediciones meteorológicas a nivel de superficie y su correspondiente header}
% \begin{tabular}{||c|c|c|c|c|c|c|c||} 
% \hline
% Record & Date/Time & 200 & Tamb(K) & Rh(\%) & Pres(mb) & Tir(K) & Rain \\ [0.5ex] 
% \hline\hline
% 25 & 10/07/14 22:08:13 & 201 & 276.032 & 10.13 & 557.48 & 191.09 & 0 \\
% \hline
% \end{tabular}
%\end{table}
%
%\newpage
%\begin{table}[h]
% \centering
% \caption{Ejemplo de mediciones de valores scalares e integrados y su correspondiente header}
% \begin{tabular}{||c|c|c|c|c|c|c||} 
% \hline
% Record & Date/Time & 300 &	Int. Vapor(cm) & Int. Liquid(mm) & Cloud Base(km) \\ [0.5ex] 
% \hline\hline
% 42	& 10/07/14 22:08:56 & 301 & 0.077 & 0 & -1 \\
% \hline
% \end{tabular}
%\end{table}
%
%\begin{table}[h]
% \centering
% \caption{Ejemplo de perfiles y sus correspondientes headers}
% \begin{tabular}{||c|c|c|c|c|c|c||} 
% \hline
% Record & Date/Time & 400 &	LV2 Processor(angle) & 0 & ... & 10 \\ [0.5ex] 
% \hline\hline
% 46 & 10/07/14 22:09:41 & 401 &	Zenith & 276.017 & ... & 259.11 \\
% \hline
% 53	& 10/07/14 22:09:43 & 402 &	Angle Scan(N) &	0.623 & ... & 0.007 \\
% \hline
% 57	& 10/07/14 22:09:43 & 403 &	Angle Scan(S) &	0.021 & ... & 0.005 \\ 
% \hline
% 61	& 10/07/14 22:09:43 & 404 &	Angle Scan(A) &	2.805 & ... & 0.336 \\
% \hline
% \end{tabular}
%\end{table}


% ============= REFERENCIAS ==============
%\newpage
%
%\section{Referencias} 
%\begin{thebibliography}{99}
%	\bibitem{nodoPag} R. Jacob Baker, \textit{CMOS. Circuit Design, Layout, and Simulation}, 2nd ed., USA: IEEE Press, 2005
%	\bibitem{Bojan} asdasdas
%	
%\end{thebibliography}


% ============= FIN DE DOCUMENTO ==============
\end{document}
